\SourcePage{103}
\SGAsection{4}{形变的形式模}
在一些性质良好的\Term{cofibered-category}{余纤维范畴}中，\Term{object}{对象}的\Term{isomorphism-class}{同构类}通常并非\Term{prorepresentable}{可预表}，但仍存在\Term{formal-versal-object}{形式万有对象}；\Term{deformation}{形变}的\Term{cofibered-category}{余纤维范畴}就是其中之一。现在回到第1节的情形与记号。因此，本节始终固定一个\Term{complete}{完备}\Term{noetherian}{诺特}\Term{local-ring}{局部环} $\Lambda$，其\Term{residue-field}{剩余域}为 $k$；$\mathcal C_\Lambda$ 表示具有相同\Term{residue-field}{剩余域} $k$ 的 Artin \Term{local-algebra}{局部 $\Lambda$-代数}的\Term{category}{范畴}。

固定域 $k$ 上的\Term{scheme}{概形} $X_0$。$X_0$ 的一个（\Term{infinitesimal-deformation}{无穷小形变}）是指一对 $(R,X)$，其中 $R\in\operatorname{ob}(\mathcal C_\Lambda)$，$X$ 是一个\Term{flat}{平坦} $R$-\Term{scheme}{概形}，并带有\Term{commutative-diagram}{交换图}
\[
\begin{tikzcd}
X_0\arrow[r,"i_X"]\arrow[d]&X\arrow[d,"\text{平坦}"]\\
\operatorname{Spec}(k)\arrow[r]&\operatorname{Spec}(R).
\end{tikzcd}
\]
\SourcePage{104}
要求诱导的\Term{morphism}{态射} $X_0\to X\times_{\operatorname{Spec}(R)}\operatorname{Spec}(k)$ 为\Term{isomorphism}{同构}。注意，因 $R$ 是 Artin \Term{local-ring}{局部环}，$i_X:X_0\to X$ 必为\Term{closed-immersion}{闭浸入}，且在底\Term{topological-space}{拓扑空间}上为\Term{homeomorphism}{同胚}。若 $(R,X)$、$(R',X')$ 分别为 $X_0$ 到 $R$、$R'$ 的\Term{deformation}{形变}，则定义\Term{arrow}{箭头} $(R',X')\to(R,X)$ 由 $\mathcal C_\Lambda$ 中一个映射 $R'\xrightarrow{\varphi}R$ 和一个\Term{scheme}{概形}\Term{morphism}{态射} $\Phi:X\to X'$ 构成，并要求下图\Term{commutative-diagram}{交换}：
\[
\begin{tikzcd}[column sep=large]
X\arrow[rr,"\Phi"]\arrow[dd]&&X'\arrow[dd]\\
&X_0\arrow[ul,"i_X"]\arrow[ur,"i_{X'}"']\arrow[d]&\\
\operatorname{Spec}(R)\arrow[rr,"\operatorname{Spec}(\varphi)"']&\operatorname{Spec}(k)\arrow[l]\arrow[r]&\operatorname{Spec}(R').
\end{tikzcd}
\]
由此得到 $X_0$ 的（\Term{infinitesimal-deformation}{无穷小形变}）\Term{category}{范畴} $\mathcal M_{X_0}$，其\Term{object}{对象}与\Term{arrow}{箭头}如上定义。此\Term{category}{范畴}配有\Term{functor}{函子} $p:\mathcal M_{X_0}\to\mathcal C_\Lambda$，由 $(R,X)\mapsto R$ 定义。现给定 $(R,X)\in\operatorname{ob}(\mathcal M_{X_0})$ 及 $\mathcal C_\Lambda$ 中一个映射 $R\xrightarrow{\varphi}S$。便得到\Term{commutative-diagram}{交换图}
\[
\begin{tikzcd}
X\times_{\operatorname{Spec}(R)}\operatorname{Spec}(S)\arrow[r,"p_1"]\arrow[d,"\text{平坦}"']&X\arrow[d,"\text{平坦}"]\\
\operatorname{Spec}(S)\arrow[r,"\operatorname{Spec}(\varphi)"']&\operatorname{Spec}(R),
\end{tikzcd}
\]
而 $\varphi$-\Term{arrow}{箭头} $(\varphi,p_1):(R,X)\to(S,X\times_{\operatorname{Spec}(R)}\operatorname{Spec}(S))$ 显然在 $\mathcal M_{X_0}$ 中是\Term{cocartesian}{余笛卡儿}的。因此 $\mathcal M_{X_0}$ 经由 $p:\mathcal M_{X_0}\to\mathcal C_\Lambda$ 成为 $\mathcal C_\Lambda$ 上的\Term{cofibered-category}{余纤维范畴}。有如下事实。

\SourcePage{105}
\textbf{4.1(a)。} 设 $(1_R,\Phi):(R,X_1)\to(R,X_2)$ 是 $\mathcal M_{X_0}$ 中的\Term{arrow}{箭头}，即 $\Phi:X_2\to X_1$ 是 $R$ 上的\Term{morphism}{态射}，且 $X_2\otimes_R k\to X_1\otimes_R k$ 为\Term{isomorphism}{同构}。因为 $R$ 为 Artin 环且 $X_i$ 是 $R$-\Term{flat}{平坦}的，立即可知 $\Phi$ 为\Term{isomorphism}{同构}，所以 $(1_R,\Phi):(R,X_1)\to(R,X_2)$ 是 $\mathcal M_{X_0}$ 中的\Term{isomorphism}{同构}。故 $\mathcal M_{X_0}$ 是 $\mathcal C_\Lambda$ 上的\Term{cofibered-groupoid}{余纤维群胚}。

\textbf{4.1(b)。} 假设 $X_0$ 是\Term{affine}{仿射}的，设 $X_0=\operatorname{Spec}(A_0)$。则对任意 $(R,X)\in\operatorname{ob}(\mathcal M_{X_0})$，$X$ 也为\Term{affine}{仿射}。因此，对任意 $R\in\operatorname{ob}(\mathcal C_\Lambda)$，\Term{category}{范畴} $\mathcal M_{X_0}(R)$ 通过取\Term{section}{截面}，\Term{canonical}{典范}\Term{equivalence-of-categories}{等价}于如下\Term{category}{范畴}：其\Term{object}{对象}为一个 $R$ 上的\Term{flat}{平坦}\Term{algebra}{代数} $A$，带有一个 $R$-\Term{algebra}{代数}映射 $j_A:A\to A_0$，诱导\Term{isomorphism}{同构} $A\otimes_R k\xrightarrow\sim A_0$；其\Term{arrow}{箭头} $\Phi:A_1\to A_2$ 为使下图\Term{commutative-diagram}{交换}的 $R$-\Term{algebra}{代数}映射：
\[
\begin{tikzcd}
A_1\arrow[r,"\Phi"]\arrow[d,"j_{A_1}"']&A_2\arrow[d,"j_{A_2}"]\\
A_0\arrow[r,equal]&A_0.
\end{tikzcd}
\]
\textbf{4.1(c)。} 考察 $\widehat{\mathcal C}_\Lambda$ 上的\Term{cofibered-groupoid}{余纤维群胚} $\widehat{\mathcal M}_{X_0}$。按定义，对任意 $R\in\operatorname{ob}(\widehat{\mathcal C}_\Lambda)$，\Term{category}{范畴} $\widehat{\mathcal M}_{X_0}(R)$ \Term{canonical}{典范}\Term{equivalence-of-categories}{等价}于\Term{formal-spectrum}{形式谱} $\operatorname{Spf}(R)$ 上的\Term{adic-formal-scheme}{进形式概形} $\mathfrak X$ 所成的\Term{category}{范畴}，要求对每个 $n\geq0$，$X_n=\mathfrak X\otimes_R R/\mathfrak m^{n+1}$ 在 $R/\mathfrak m^{n+1}$ 上\Term{flat}{平坦}（若 $X_0$ 为\Term{noetherian}{诺特}，这\Term{equivalence}{等价}于 $\mathfrak X$ 在 $\operatorname{Spf}(R)$ 上\Term{flat}{平坦}）。故若 $X_0$ 为\Term{noetherian}{诺特}，则 $\widehat{\mathcal M}_{X_0}$ \Term{canonical}{典范}地 $\widehat{\mathcal C}_\Lambda$-\Term{equivalence-of-categories}{等价}于\Term{formal-scheme}{形式概形}意义下 $X_0$ 在 $\widehat{\mathcal C}_\Lambda$ 的\Term{object}{对象}上的\Term{deformation}{形变}\Term{category}{范畴}。

\textbf{4.1(d)。} 设 $(R,X)\in\operatorname{ob}(\mathcal M_{X_0})$，即有\Term{commutative-diagram}{交换图}
\SourcePage{106}
\[
\begin{tikzcd}
X_0\arrow[r,"i_X"]\arrow[d]&X\arrow[d,"\text{平坦}"]\\
\operatorname{Spec}(k)\arrow[r]&\operatorname{Spec}(R),
\end{tikzcd}
\]
使 $X_0\to X\times_{\operatorname{Spec}(R)}\operatorname{Spec}(k)$ 为\Term{isomorphism}{同构}。固定 $X_0$ 的\Term{open-subscheme}{开子概形} $U_0\subset X_0$。因 $i_X:X_0\to X$ 是在底\Term{topological-space}{拓扑空间}上为\Term{homeomorphism}{同胚}的\Term{closed-immersion}{闭浸入}，得到\Term{commutative-diagram}{交换图}
\[
\begin{tikzcd}
U_0\arrow[r,"i_X|_{U_0}"]\arrow[d]&X|_{U_0}\arrow[d,"\text{平坦}"]\\
\operatorname{Spec}(k)\arrow[r]&\operatorname{Spec}(R),
\end{tikzcd}
\]
其中 $U_0\to(X|_{U_0})\times_{\operatorname{Spec}(R)}\operatorname{Spec}(k)=(X\otimes_R k)|_{U_0}$ 为\Term{isomorphism}{同构}，故 $(R,X|_{U_0})$ 为 $U_0$ 到 $R$ 的\Term{deformation}{形变}。此外，对 $\mathcal M_{X_0}$ 中任意\Term{arrow}{箭头} $(\varphi,\Phi):(R,X)\to(R',X')$，$(\varphi,\Phi|_{U_0}):(R,X|_{U_0})\to(R',X'|_{U_0})$ 是 $\mathcal M_{U_0}$ 中的\Term{arrow}{箭头}。因此有 $\mathcal C_\Lambda$ 上的\Term{restriction-functor}{限制函子} $\mathcal M_{X_0}\to\mathcal M_{U_0}$。还可考虑在一点处的\Term{localization}{局部化}和\Term{formalization}{形式化}。具体说，固定 $X_0$ 中一点 $x_0$。若 $(R,X)\in\operatorname{ob}(\mathcal M_{X_0})$，则 $\operatorname{Spec}(\mathcal O_{X,x_0})$ 是 $\operatorname{Spec}(\mathcal O_{X_0,x_0})$ 到 $R$ 的\Term{deformation}{形变}，$\operatorname{Spec}(\widehat{\mathcal O}_{X,x_0})$ 是 $\operatorname{Spec}(\widehat{\mathcal O}_{X_0,x_0})$ 在 $R$ 上的\Term{deformation}{形变}，这里 $\widehat{\phantom A}$ 表示对由\Term{maximal-ideal}{极大理想}的幂给出的\Term{linear}{线性}\Term{topology}{拓扑}取\Term{completion}{完备化}。由此得到\Term{functor}{函子}
\[
\mathcal M_{X_0}\longrightarrow\mathcal M_{\operatorname{Spec}(\mathcal O_{X_0,x_0})}\longrightarrow\mathcal M_{\operatorname{Spec}(\widehat{\mathcal O}_{X_0,x_0})}.
\]
首先证明 $\mathcal M_{X_0}$ 是 $\mathcal C_\Lambda$ 上的\Term{homogeneous}{齐性}\Term{cofibered-category}{余纤维范畴}。

\begin{SGAresult}{\Term{lemma}{引理}}{4.2}
考察 $\mathcal C_\Lambda$ 中的图
\[
\begin{tikzcd}
&R'\arrow[d]\\
R\arrow[r,two heads]&\overline R,
\end{tikzcd}
\]
设 $E$、$E'$ 分别是 $R$、$R'$ 上的\Term{flat}{平坦}\Term{module}{模}。则对任意 $R'$-\Term{homomorphism}{同态} $E'\to\overline E=E\otimes_R\overline R$，\Term{fiber-product}{纤维积} $E\times_{\overline E}E'$ 是\Term{flat}{平坦} $R\times_{\overline R}R'$-\Term{module}{模}。
\end{SGAresult}
\SourcePage{107}
\begin{proof}
设 $0\to I\to R\to\overline R\to0$ \Term{exact}{正合}，则 $0\to I\to R\times_{\overline R}R'\to R'\to0$ \Term{exact}{正合}。因为 $E$ 是 $R$-\Term{flat}{平坦}的，$0\to I\otimes_R E\to E\to\overline E\to0$ \Term{exact}{正合}，于是得到\Term{exact-sequence}{正合列}
\[
0\longrightarrow I\otimes_R E\longrightarrow E\times_{\overline E}E'\longrightarrow E'\longrightarrow0.\tag{*}
\]
令 $S=R\times_{\overline R}R'$。因 $I$ 为\Term{nilpotent}{幂零}，$E\times_{\overline E}E'$ 作为 $S$-\Term{module}{模}\Term{flat}{平坦}，当且仅当 $(E\times_{\overline E}E')\otimes_S R'$ 是\Term{flat}{平坦} $R'$-\Term{module}{模}，且 $\operatorname{Tor}_1^S(E\times_{\overline E}E',R')=0$。但注意到
\[
I\otimes_R E=I\otimes_S(E\times_{\overline E}E'),\qquad
E'=R'\otimes_S(E\times_{\overline E}E'),
\]
故\Term{exact-sequence}{正合列} (*) 可改写为
\[
0\longrightarrow I\otimes_S(E\times_{\overline E}E')\longrightarrow E\times_{\overline E}E'\longrightarrow R'\otimes_S(E\times_{\overline E}E')\longrightarrow0.
\]
这就是说 $\operatorname{Tor}_1^S(R',E\times_{\overline E}E')=0$，且 $R'\otimes_S(E\times_{\overline E}E')=E'$ 是\Term{flat}{平坦} $R'$-\Term{module}{模}，故 $E\times_{\overline E}E'$ 是 $S$-\Term{flat}{平坦}的。
\end{proof}

\begin{SGAresult}{推论}{4.3}
给定 $\mathcal M_{X_0}$ 中的图
\[
\begin{tikzcd}
(R,X)\arrow[dr,"{(\varphi,\Phi)}"']&&(R',X')\arrow[dl,"{(\varphi',\Phi')}"]\\
&(\overline R,\overline X),&
\end{tikzcd}
\]
若 $\varphi:R\to\overline R$ 为\Term{surjectivity}{满射}，则\Term{amalgamated-sum}{融合和} $X\amalg_{\overline X}X'$（在\Term{scheme}{概形}意义下）是 $R\times_{\overline R}R'$ 上的\Term{flat}{平坦}\Term{scheme}{概形}。因此，$\mathcal M_{X_0}$ 是 $\mathcal C_\Lambda$ 上的\Term{homogeneous}{齐性}\Term{cofibered-groupoid}{余纤维群胚}。
\end{SGAresult}
\begin{proof}
对每个\Term{affine-open}{仿射开集} $U_0\subset X_0$，由定义有
\[
(X\amalg_{\overline X}X')|_{U_0}\simeq
\operatorname{Spec}\bigl(\Gamma(U_0,X)\times_{\Gamma(U_0,\overline X)}\Gamma(U_0,X')\bigr).
\]
因此，若 $\varphi:R\to\overline R$ 为\Term{surjectivity}{满射}，由4.2知 $X\amalg_{\overline X}X'$ 是 $R\times_{\overline R}R'$ 上的\Term{flat}{平坦}\Term{scheme}{概形}，且\Term{morphism}{态射} $i_X\amalg_{i_{\overline X}}i_{X'}:X_0\to X\amalg_{\overline X}X'$ 诱导\Term{isomorphism}{同构} $X_0\xrightarrow\sim(X\amalg_{\overline X}X')\otimes_S k$，其中 $S=R\times_{\overline R}R'$。所以 $(R\times_{\overline R}R',X\amalg_{\overline X}X')\in\operatorname{ob}(\mathcal M_{X_0})$，而显然在\Term{category}{范畴} $\mathcal M_{X_0}$ 中有
\[
(R\times_{\overline R}R',X\amalg_{\overline X}X')
=(R,X)\times_{(\overline R,\overline X)}(R',X'),
\]
\SourcePage{108}
这正是 $\mathcal M_{X_0}$ 在 $\mathcal C_\Lambda$ 上\Term{homogeneous}{齐性}。
\end{proof}

\begin{SGAresult}{\Term{lemma}{引理}}{4.4}
对域 $k$ 上的任意\Term{scheme}{概形} $X_0$，有\Term{canonical}{典范}\Term{exact-sequence}{正合列}
\[
0\longrightarrow H^1(X_0,\mathcal D^0_{X_0})\longrightarrow
[\mathcal M_{X_0}(k[\epsilon])]\longrightarrow H^0(X_0,\mathcal D^1_{X_0}),
\]
其中 $\mathcal D^i_{X_0}$ 如第3节所定义。
\end{SGAresult}
\begin{proof}
（特殊情形。）假设 $X_0$ 为\Term{affine}{仿射}，令 $\Gamma(X_0,\mathcal O_{X_0})=A_0$。如4.1(b)中所述，\Term{category}{范畴} $\mathcal M_{X_0}(k[\epsilon])$ \Term{canonical}{典范}\Term{equivalence-of-categories}{等价}于\Term{category}{范畴} $\mathcal E$，其中：

$\operatorname{ob}(\mathcal E)$：$\mathcal E$ 的一个\Term{object}{对象}是 $k[\epsilon]$ 上的\Term{flat}{平坦}\Term{algebra}{代数} $A$，带有 $k[\epsilon]$-\Term{algebra}{代数}映射 $j_A:A\to A_0$，诱导\Term{isomorphism}{同构}
\[
k\otimes_{k[\epsilon]}j_A:k\otimes_{k[\epsilon]}A\xrightarrow\sim A_0.
\]
也可将 $\mathcal E$ 的\Term{object}{对象}描述为一个 $k$-\Term{algebra}{代数} $A$，带有\Term{exact-sequence}{正合列}
\[
0\longrightarrow A_0\xrightarrow{\epsilon}A\xrightarrow{j_A}A_0\longrightarrow0,
\]
其中 $\epsilon(A_0)^2=0$，这里将 $\epsilon(A_0)$ 看作 $A$ 的\Term{ideal}{理想}。

$\operatorname{Fl}(\mathcal E)$：按照 $\mathcal E$ 的\Term{object}{对象}的第一种描述，$\mathcal E$ 中一个\Term{arrow}{箭头} $A_1\to A_2$ 是 $k[\epsilon]$-\Term{algebra}{代数}映射 $\Phi:A_1\to A_2$，满足 $j_{A_1}=j_{A_2}\circ\Phi$。所以，按\Term{object}{对象}的另一种描述，$\mathcal E$ 中一个\Term{arrow}{箭头} $A_1\to A_2$ 是使下图\Term{commutative-diagram}{交换}的 $k$-\Term{algebra}{代数}映射 $\Phi:A_1\to A_2$：
\[
\begin{tikzcd}
0\arrow[r]&A_0\arrow[r,"\epsilon"]\arrow[d,equal]&A_1\arrow[r,"j_{A_1}"]\arrow[d,"\Phi"]&A_0\arrow[r]\arrow[d,equal]&0\\
0\arrow[r]&A_0\arrow[r,"\epsilon"]&A_2\arrow[r,"j_{A_2}"]&A_0\arrow[r]&0.
\end{tikzcd}
\]
因此，$[\mathcal M_{X_0}(k[\epsilon])]$ 经由 $\Gamma(X_0,-)$ \Term{canonical}{典范}\Term{isomorphism}{同构}于 $A_0$ 被 $A_0$ \Term{group-extension}{扩张}的 $k$-\Term{algebra}{代数}\Term{extension}{扩张}的\Term{isomorphism-class}{同构类}之集。换言之，经由 $\Gamma(X_0,-)$ 有\Term{canonical}{典范}\Term{bijection}{双射}
\[
[\mathcal M_{X_0}(k[\epsilon])]\xrightarrow\sim H^0(X_0,\mathcal D^1_{X_0}),
\]
且显然保持两边的\Term{vector-space}{向量空间}结构。

\SourcePage{109}
（\Term{generic}{一般}情形。）对每个\Term{affine-open}{仿射开集} $U_0\subset X_0$，由复合映射
\[
[\mathcal M_{X_0}(k[\epsilon])]\longrightarrow[\mathcal M_{U_0}(k[\epsilon])]
\xrightarrow\sim H^0(U_0,\mathcal D^1_{X_0})
\]
得到\Term{canonical-map}{典范映射} $[\mathcal M_{X_0}(k[\epsilon])]\to H^0(U_0,\mathcal D^1_{X_0})$，从而得到\Term{canonical-map}{典范映射}
\[
[\mathcal M_{X_0}(k[\epsilon])]\longrightarrow H^0(X_0,\mathcal D^1_{X_0}).
\]
其\Term{kernel}{核}是 $X_0$ 到 $k[\epsilon]$ 的\Term{locally-trivial-deformation}{局部平凡形变}的\Term{isomorphism-class}{同构类}之集，即那些\Term{local-scheme}{局部}（在 $X$ 上）\Term{isomorphic}{同构}于 $X_0\otimes_k k[\epsilon]$ 的 $X_0$ 到 $k[\epsilon]$ 的\Term{deformation}{形变}的\Term{isomorphism-class}{同构类}之集。因为
\[
\mathcal D^0_{X_0}=\operatorname{Der}_k(\mathcal O_{X_0},\mathcal O_{X_0})
=\operatorname{Aut}(X_0\otimes_k k[\epsilon]\mid X_0)
\]
是 $X_0\otimes_k k[\epsilon]$ 的在\Term{subscheme}{子概形} $X_0$ 上诱导恒等的\Term{automorphism}{自同构}的芽所成的\Term{sheaf}{层}，上述\Term{set}{集合}\Term{canonical}{典范}\Term{isomorphism}{同构}于 $H^1(X_0,\mathcal D^0_{X_0})$，便得到所需的\Term{exact-sequence}{正合列}。
\end{proof}

\begin{SGAresult}{\Term{theorem}{定理}}{4.5}
设 $X_0$ 是 $k$ 上的\Term{finite-type}{有限型}\Term{scheme}{概形}，满足以下条件之一：(i) $X_0$ 在 $k$ 上\Term{proper}{固有}；(ii) $X_0$ 为\Term{affine}{仿射}，且
\[
X_0^{\mathrm{sing}}=\{x\in X_0\mid X_0\text{ 在点 }x\text{ 处不在 }k\text{ 上光滑}\}
\]
为\Term{finite}{有限}集。则：

(a) $\mathcal M_{X_0}$ 在 $\mathcal C_\Lambda$ 上具有\Term{formal-versal-object}{形式万有对象}（称为 $X_0$ 的一个\Term{formal-versal-deformation}{形式万有形变}）。若 $(R,\mathfrak X)$ 是 $X_0$ 的\Term{formal-versal-deformation}{形式万有形变}，则对任意 $(S,\mathfrak Y)\in\operatorname{ob}(\widehat{\mathcal M}_{X_0})$，存在 $\widehat{\mathcal M}_{X_0}$ 中的映射 $(\varphi,\Phi):(R,\mathfrak X)\to(S,\mathfrak Y)$，使下图为\Term{cartesian}{笛卡儿}图：
\[
\begin{tikzcd}
\mathfrak Y\arrow[r,"\Phi"]\arrow[d]&\mathfrak X\arrow[d]\\
\operatorname{Spf}(S)\arrow[r,"\operatorname{Spf}(\varphi)"']&\operatorname{Spf}(R).
\end{tikzcd}
\]
此外，$X_0$ 的一个\Term{formal-versal-deformation}{形式万有形变} $(R,\mathfrak X)$ \Term{prorepresentable}{预表}\Term{functor}{函子} $[\mathcal M_{X_0}]$，当且仅当对 $\mathcal C_\Lambda$ 中每个\Term{surjectivity}{满射} $R'\to R$ 及 $X_0$ 到 $R'$ 的\Term{deformation}{形变} $(R',X')$，\Term{homomorphism}{同态}
\[
\operatorname{Aut}_{R'}(X'\mid X_0)\longrightarrow
\operatorname{Aut}_R(X'\otimes_{R'}R\mid X_0)
\]
是\Term{surjectivity}{满射}。

(b) 若 $X_0$ 在 $k$ 上\Term{proper}{固有}，则 $\mathcal M_{X_0}$ 在 $\mathcal C_\Lambda$ 上\Term{smoothly-pseudo-representable}{光滑伪可表}。
\end{SGAresult}
\SourcePage{110}
\begin{proof}
(a) 由4.3，$\mathcal M_{X_0}$ 为 $\mathcal C_\Lambda$ 上的\Term{homogeneous}{齐性}\Term{cofibered-groupoid}{余纤维群胚}，故由1.11及2.9(b)，只需证明 $[\mathcal M_{X_0}(k[\epsilon])]$ 是 $k$ 上的\Term{finite-dimensional}{有限维}\Term{vector-space}{向量空间}。

考察\Term{exact-sequence}{正合列}
\[
0\longrightarrow H^1(X_0,\mathcal D^0_{X_0})\longrightarrow
[\mathcal M_{X_0}(k[\epsilon])]\longrightarrow H^0(X_0,\mathcal D^1_{X_0}).\tag{*}
\]
若 $X_0$ 在 $k$ 上\Term{proper}{固有}，则 $H^1(X_0,\mathcal D^0_{X_0})$ 与 $H^0(X_0,\mathcal D^1_{X_0})$ 都\Term{finite-dimensional}{有限维}，因为由2.13(a)，$\mathcal D^i_{X_0}$ 对所有 $i\geq0$ 都\Term{coherent}{凝聚}。现假设 $X_0$ 为\Term{affine}{仿射}且 $X_0^{\mathrm{sing}}$ 为\Term{finite}{有限}集。由2.13(a)，$\mathcal D^1_{X_0}$ 是 $X_0$ 上的\Term{coherent}{凝聚}\Term{sheaf}{层}且 $\operatorname{Supp}\mathcal D^1_{X_0}\subset X_0^{\mathrm{sing}}$，故 $H^0(X_0,\mathcal D^1_{X_0})$ 是 $k$ 上的\Term{finite-dimensional}{有限维}\Term{vector-space}{向量空间}；又因 $X_0$ 为\Term{affine}{仿射}，$H^1(X_0,\mathcal D^0_{X_0})=0$。因此两种情形下都由\Term{exact-sequence}{正合列} (*) 得到 $[\mathcal M_{X_0}(k[\epsilon])]$ 是 $k$ 上的\Term{finite-dimensional}{有限维}\Term{vector-space}{向量空间}。(a)的第二个断言由2.7得到。

(b) 若 $X_0$ \Term{proper}{固有}，则
\[
\operatorname{Aut}(X_0\otimes_k k[\epsilon]\mid X_0)=H^0(X_0,\mathcal D^0_{X_0})
\]
为\Term{finite-dimensional}{有限维}\Term{vector-space}{向量空间}，故断言由2.13得到。
\end{proof}

\textbf{注记 4.6(a)。} 设 $X_0$ 为\Term{affine}{仿射}。则 $X_0$ 的\Term{formal-versal-deformation}{形式万有形变}存在，当且仅当 $\mathcal D^1_{X_0}$ 的\Term{support}{支集}由\Term{finite}{有限}个点组成。但这不意味着 $X_0$ 的所有非\Term{smooth}{光滑}点都孤立：有些 $k$ 上的\Term{affine}{仿射}\Term{scheme}{概形} $X_0$ 具有非孤立的非\Term{smooth}{光滑}点，而 $\operatorname{Supp}\mathcal D^1_{X_0}$ 却为\Term{finite}{有限}集。

\textbf{注记 4.6(b)。} 此后，$X_0$ 的\Term{deformation}{形变}指 $\widehat{\mathcal M}_{X_0}$ 中的一个\Term{object}{对象}。可能发生混淆时，将 $\mathcal M_{X_0}$ 中的\Term{object}{对象}称为 $X_0$ 的\Term{infinitesimal-deformation}{无穷小形变}。

\textbf{注记 4.6(c)。} 设 $X_0$ 是 $k$ 上的\Term{scheme}{概形}。我们固定了一个\Term{residue-field}{剩余域}为 $k$ 的\Term{complete}{完备}\Term{noetherian}{诺特}\Term{local-ring}{局部环} $\Lambda$，并考虑 $X_0$ 在\Term{category}{范畴} $\mathcal C_\Lambda$ 上的\Term{deformation}{形变}。
\SourcePage{111}
（实际中通常取 $\Lambda=k$，或在 $k$ 为正\Term{characteristic}{特征}\Term{perfect-field}{完美域}时取 $\Lambda=W(k)$。）$\mathcal C_k\subset\mathcal C_\Lambda$ 是一个\Term{fully-faithful}{全忠实}\Term{embedding}{嵌入}，而 $\mathcal M_{X_0}|_{\mathcal C_k}$ 就是 $X_0$ 在 $S\in\operatorname{ob}(\mathcal C_k)$ 上的\Term{deformation}{形变}所成的\Term{category}{范畴}。因此，若 $(R,\mathfrak X)$ 是 $X_0$ 在 $\mathcal C_\Lambda$ 上的\Term{formal-versal-deformation}{形式万有形变}，则 $(k\otimes_\Lambda R,k\otimes_\Lambda\mathfrak X)$ 是 $X_0$ 在 $\mathcal C_k$ 上的\Term{formal-versal-deformation}{形式万有形变}。例如，由此可知
\[
\dim_{\mathcal C_\Lambda}\mathcal M_{X_0}=\dim k\otimes_\Lambda R
\]
与 $\Lambda$ 的选择无关，只依赖于 $X_0$。

上述\Term{theorem}{定理}可以向许多方向推广。例如考虑一对 $Y_0\to X_0$，其中 $X_0$ 为 $k$ 上的\Term{finite-type}{有限型}\Term{scheme}{概形}，$Y_0$ 为 $X_0$ 的一个固定\Term{closed-subscheme}{闭子概形}。对每个 $R\in\operatorname{ob}(\mathcal C_\Lambda)$，$(X_0,Y_0)$ 到 $R$ 的一个\Term{deformation}{形变}指\Term{commutative-diagram}{交换图}
\[
\begin{tikzcd}[column sep=large,row sep=large]
Y_0\arrow[r]\arrow[d]\arrow[dd,bend right=45]&Y\arrow[d]\arrow[dd,bend left=45,"\text{平坦}"]\\
X_0\arrow[r]\arrow[d]&X\arrow[d,"\text{平坦}"]\\
\operatorname{Spec}(k)\arrow[r]&\operatorname{Spec}(R).
\end{tikzcd}
\]
要求 $(X_0,Y_0)\to(X\times_{\operatorname{Spec}(R)}\operatorname{Spec}(k),Y\times_{\operatorname{Spec}(R)}\operatorname{Spec}(k))$ 为\Term{isomorphism}{同构}。以显然的方式定义 $(X_0,Y_0)$ 的\Term{deformation}{形变}之间的\Term{morphism}{态射} $(R';X',Y')\to(R;X,Y)$，便得到 $\mathcal C_\Lambda$ 上的\Term{cofibered-groupoid}{余纤维群胚} $\mathcal M_{X_0,Y_0}$。（这是\Term{deformation}{形变}的模问题中经常用于\Term{rigidification}{刚化}\Term{automorphism}{自同构}的过程。）以4.2为基础，使用与4.3相同的论证，立即看出 $\mathcal M_{X_0,Y_0}$ 为 $\mathcal C_\Lambda$ 上的\Term{homogeneous}{齐性}\Term{cofibered-groupoid}{余纤维群胚}。

\begin{SGAresult}{\Term{lemma}{引理}}{4.7}
序列
\[
0\longrightarrow H^0\bigl(Y_0,\mathcal Hom_{\mathcal O_{Y_0}}(I/I^2,\mathcal O_{Y_0})\bigr)
\longrightarrow[\mathcal M_{X_0,Y_0}(k[\epsilon])]
\longrightarrow[\mathcal M_{X_0}(k[\epsilon])]
\]
\Term{exact}{正合}，其中 $I$ 是 $\mathcal O_{X_0}$ 中定义 $Y_0$ 的\Term{ideal-sheaf}{理想层}，即
\SourcePage{112}
$0\to I\to\mathcal O_{X_0}\to\mathcal O_{Y_0}\to0$ \Term{exact}{正合}。
\end{SGAresult}
\begin{proof}
设 $\mathcal B$ 为 $\mathcal M_{X_0,Y_0}(k[\epsilon])$ 的\Term{full-subcategory}{全子范畴}，\Term{object}{对象}形如 $(k[\epsilon];X\times_{\operatorname{Spec}(k)}\operatorname{Spec}(k[\epsilon]),Y)$。于是
\[
[\mathcal B]=\operatorname{Ker}\bigl([\mathcal M_{X_0,Y_0}(k[\epsilon])]
\longrightarrow[\mathcal M_{X_0}(k[\epsilon])]\bigr).
\]
\Term{category}{范畴} $\mathcal B$ 是 $Y_0$ 到 $k[\epsilon]$ 的“在 $X_0\times_{\operatorname{Spec}(k)}\operatorname{Spec}(k[\epsilon])$ 内”的\Term{deformation}{形变}的\Term{groupoid}{群胚}，即 $\mathcal B$ 的一个\Term{object}{对象}是\Term{commutative-diagram}{交换图}
\[
\begin{tikzcd}
Y\arrow[rr,"i_Y"]\arrow[dr,"\text{平坦}"']&&X_0\times_{\operatorname{Spec}(k)}\operatorname{Spec}(k[\epsilon])\arrow[dl]\\
&\operatorname{Spec}(k[\epsilon]),&
\end{tikzcd}
\]
满足 $\operatorname{Spec}(k)\times_{\operatorname{Spec}(k[\epsilon])}i_Y=i_{Y_0}$；一个\Term{arrow}{箭头} $(Y_1,i_{Y_1})\to(Y_2,i_{Y_2})$ 是 $k[\epsilon]$-\Term{morphism}{态射} $\Phi:Y_1\to Y_2$（必为\Term{isomorphism}{同构}），满足
\[
\operatorname{Spec}(k)\times_{\operatorname{Spec}(k[\epsilon])}\Phi=\operatorname{id}_{Y_0},\qquad
 i_{Y_2}\circ\Phi=i_{Y_1}.
\]
也可将 $\mathcal B$ 的\Term{object}{对象}描述为一对 $(Y,j_Y)$，其中 $Y$ 是 $Y_0$ 到 $k[\epsilon]$ 的\Term{deformation}{形变}，$j_Y:Y\to X_0$ 是 $k$ 上的\Term{morphism}{态射}，使复合\Term{morphism}{态射} $Y_0\to Y\xrightarrow{j_Y}X_0$ 为 $i_{Y_0}$；一个\Term{arrow}{箭头} $(Y_1,j_{Y_1})\to(Y_2,j_{Y_2})$ 是\Term{morphism}{态射} $\Phi:Y_1\to Y_2$，满足 $k\otimes_{k[\epsilon]}\Phi=\operatorname{id}_{Y_0}$ 及 $j_{Y_2}\circ\Phi=j_{Y_1}$。故对应
\[
(Y,j_Y)\longmapsto
\begin{tikzcd}[column sep=small]
0\arrow[r]&I\arrow[r]\arrow[d]&\mathcal O_{X_0}\arrow[r]\arrow[d,"j_Y"]&\mathcal O_{Y_0}\arrow[r]\arrow[d,equal]&0\\
0\arrow[r]&\mathcal O_{Y_0}\arrow[r,"\epsilon"]&\mathcal O_Y\arrow[r]&\mathcal O_{Y_0}\arrow[r]&0
\end{tikzcd}
\]
给出\Term{isomorphism}{同构}
\[
\begin{aligned}
[\mathcal B]&\xrightarrow{\ j_Y|_I\ }\operatorname{Hom}_{\mathcal O_{X_0}}(I,\mathcal O_{Y_0})
=\operatorname{Hom}_{\mathcal O_{Y_0}}(I/I^2,\mathcal O_{Y_0})\\
&=H^0\bigl(Y_0,\mathcal Hom_{\mathcal O_{Y_0}}(I/I^2,\mathcal O_{Y_0})\bigr).
\end{aligned}
\]
\end{proof}

\SourcePage{113}
\textbf{注记 4.8。} 蕴含上述\Term{lemma}{引理}的一个更\Term{generic}{一般}的公式见 SGA~3 第 III 讲4.4。还可注意如下事实：对任意 $S$-\Term{scheme}{概形} $X$ 及\Term{quasi-coherent}{拟凝聚} $X$-\Term{module}{模} $E$，可对所有 $i\geq0$ 定义\Term{quasi-coherent}{拟凝聚} $X$-\Term{module}{模} $\mathcal D^i_{X\mid S}(E)$，推广 $S$ 为\Term{affine}{仿射}的情形。照例简记 $\mathcal D^i_{X\mid S}(\mathcal O_X)=\mathcal D^i_{X\mid S}$。于是
\[
H^0\bigl(Y_0,\mathcal Hom_{\mathcal O_{Y_0}}(I/I^2,\mathcal O_{Y_0})\bigr)
=H^0(Y_0,\mathcal D^1_{Y_0\mid X}),
\]
其中 $0\to I\to\mathcal O_{X_0}\to\mathcal O_{Y_0}\to0$ \Term{exact}{正合}。

\begin{SGAresult}{\Term{theorem}{定理}}{4.9}
设 $X_0$ 为 $k$ 上的\Term{finite-type}{有限型}\Term{scheme}{概形}，$Y_0$ 为 $X_0$ 的\Term{closed-subscheme}{闭子概形}。假设满足以下条件之一：(i) $X_0$ 在 $k$ 上\Term{proper}{固有}；(ii) $X_0$ 为\Term{affine}{仿射}且在\Term{finite}{有限}个点以外在 $k$ 上\Term{smooth}{光滑}，而 $Y_0$ 在 $k$ 上\Term{proper}{固有}。则有：

(a) $\mathcal M_{X_0,Y_0}$ 在 $\mathcal C_\Lambda$ 上有\Term{formal-versal-deformation}{形式万有形变}，且它\Term{prorepresentable}{预表}\Term{functor}{函子} $[\mathcal M_{X_0,Y_0}]$，当且仅当对 $\mathcal C_\Lambda$ 中每个\Term{surjectivity}{满射} $R'\to R$ 和一个\Term{deformation}{形变} $(R';X',Y')$，映射
\[
\operatorname{Aut}_{R'}((X',Y')\mid X_0)\longrightarrow
\operatorname{Aut}_R((X'\otimes_{R'}R,Y'\otimes_{R'}R)\mid X_0)
\]
是\Term{surjectivity}{满射}。

(b) 若还满足
\[
H^0\left(X_0,
\mathcal D^0_{Y_0\mid k}\mathbin{\times}_{\mathcal D^0_{X_0\mid k}(\mathcal O_{Y_0})}
\mathcal D^0_{X_0\mid k}\right)
\]
\Term{finite-dimensional}{有限维}，则 $\mathcal M_{X_0,Y_0}$ \Term{smoothly-pseudo-representable}{光滑伪可表}。
\end{SGAresult}
\begin{proof}
由4.4、4.6以及1.11、2.7、2.13立即得到。
\end{proof}

现在考虑\Term{scheme}{概形}的\Term{deformation-theory}{形变理论}中向\Term{localization}{局部化}及\Term{completion}{完备化}的过渡。先说明一个记号：设 $X_0$ 为 $k$ 上的\Term{finite-type}{有限型}\Term{scheme}{概形}。给定一个固定点 $x_0\in X_0$，有自然\Term{functor}{函子}
\[
\mathcal M_{X_0}\longrightarrow\mathcal M_{\operatorname{Spec}(\mathcal O_{X_0,x_0})}
\longrightarrow\mathcal M_{\operatorname{Spec}(\widehat{\mathcal O}_{X_0,x_0})},
\]
\SourcePage{114}
从而得到\Term{functor}{函子}
\[
\widehat{\mathcal M}_{X_0}\longrightarrow\widehat{\mathcal M}_{\operatorname{Spec}(\mathcal O_{X_0,x_0})}
\longrightarrow\widehat{\mathcal M}_{\operatorname{Spec}(\widehat{\mathcal O}_{X_0,x_0})}.
\]
如4.1(c)中所述，$\widehat{\mathcal M}_{X_0}$ \Term{canonical}{典范}\Term{equivalence-of-categories}{等价}于如下\Term{category}{范畴}：一个\Term{object}{对象}为一对 $(R,\mathfrak X)$，其中 $R\in\operatorname{ob}(\widehat{\mathcal C}_\Lambda)$，$\mathfrak X$ 为一个\Term{flat}{平坦}的 $R$-\Term{adic-formal-scheme}{进形式概形}，并带有一个固定\Term{isomorphism}{同构} $\mathfrak X_0\xrightarrow\sim X_0$。此后，在不会混淆时，将 $(R,\mathfrak X)$ 在
\[
\widehat{\mathcal M}_{X_0}\longrightarrow\widehat{\mathcal M}_{\operatorname{Spec}(\mathcal O_{X_0,x_0})},\qquad
\widehat{\mathcal M}_{X_0}\longrightarrow\widehat{\mathcal M}_{\operatorname{Spec}(\widehat{\mathcal O}_{X_0,x_0})}
\]
下的像分别记为 $(R,\mathfrak X_{(x_0)})$ 及 $(R,\mathfrak X_{(\widehat{x}_0)})$。因此，若
\[
\mathfrak X=\varinjlim_\nu\mathfrak X_\nu,\qquad
\mathfrak X_\nu=\mathfrak X\times_{\operatorname{Spf}(R)}\operatorname{Spec}(R/\mathfrak m^{\nu+1}),
\]
则
\[
\mathfrak X_{(x_0)}=\varinjlim_\nu\operatorname{Spec}(\mathcal O_{\mathfrak X_\nu,x_0}),\qquad
\mathfrak X_{(\widehat{x}_0)}=\varinjlim_\nu\operatorname{Spec}(\widehat{\mathcal O}_{\mathfrak X_\nu,x_0}).
\]

\begin{SGAresult}{\Term{theorem}{定理}}{4.10}
设 $X_0$ 为 $k$ 上的\Term{finite-type}{有限型}\Term{affine}{仿射}\Term{scheme}{概形}。

(a) 若 $X_0$ \Term{local-scheme}{局部}是\Term{complete-intersection}{完全交}，则 $\mathcal M_{X_0}$ 在 $\mathcal C_\Lambda$ 上\Term{smooth}{光滑}。

(b) 假设对某个固定\Term{closed-point}{闭点} $x\in X_0$，$X_0-\{x\}$ 在 $k$ 上\Term{smooth}{光滑}。则
\[
\mathcal M_{X_0}\longrightarrow\mathcal M_{\operatorname{Spec}(\mathcal O_{X_0,x})}
\longrightarrow\mathcal M_{\operatorname{Spec}(\widehat{\mathcal O}_{X_0,x})}
\]
两者均为\Term{minimally-smooth}{极小光滑}\Term{functor}{函子}。特别地，若 $(R,\mathfrak X)$ 为 $X_0$ 的\Term{formal-versal-deformation}{形式万有形变}，则 $(R,\mathfrak X_{(x)})$ 及 $(R,\mathfrak X_{(\widehat x)})$ 分别为 $\operatorname{Spec}(\mathcal O_{X_0,x})$ 及 $\operatorname{Spec}(\widehat{\mathcal O}_{X_0,x})$ 的\Term{formal-versal-deformation}{形式万有形变}。
\end{SGAresult}
\begin{proof}
本证明回到第3节定义的交换\Term{algebra}{代数}\Term{cohomology}{上同调}的记号。令 $X_0=\operatorname{Spec}(A_0)$。

\SourcePage{115}
(a) 只需证明 $P:[\mathcal M_{X_0}]\to\mathcal C_\Lambda$ \Term{smooth}{光滑}。设 $(R,\operatorname{Spec}(A))$ 是 $\operatorname{Spec}(A_0)$ 到 $R$ 的\Term{deformation}{形变}，其中 $R\in\operatorname{ob}(\mathcal C_\Lambda)$。给定 $\mathcal C_\Lambda$ 中的\Term{small-extension}{小扩张} $R'\to R$（即 $0\to k\to R'\to R\to0$ \Term{exact}{正合}），须证明 $(R,\operatorname{Spec}(A))$ \Term{liftable}{可提升}到 $R'$。将系数为 $A_0$ 的 $D^i$ 的\Term{exact-sequence}{正合列}应用于环\Term{homomorphism}{同态} $R'\to R\to A$，得到\Term{exact-sequence}{正合列}
\[
\cdots\longrightarrow D^1(A\mid R,A_0)\longrightarrow D^1(A\mid R',A_0)
\longrightarrow D^1(R\mid R',k)\otimes_k A_0
\longrightarrow D^2(A\mid R,A_0)\longrightarrow\cdots.
\]
因 $R'\to R$ 是\Term{small-extension}{小扩张}，$D^1(R'\mid R,k)$ 是 $k$ 上的一维\Term{vector-space}{向量空间}，由\Term{extension}{扩张} $0\to k\to R'\to R\to0$ 所表示的元素（记作 $[R']$）生成。另一方面，$A$ 是 $R$-\Term{flat}{平坦}的，故对所有 $i>0$ 有
\[
D^i(A\mid R,A_0)=D^i(A_0\mid k,A_0),
\]
特别地，由3.12(11)，$D^2(A\mid R,A_0)=0$，因为 $A_0$ \Term{local-scheme}{局部}是\Term{complete-intersection}{完全交}。于是得到\Term{exact-sequence}{正合列}
\[
\longrightarrow D^1(A_0\mid k,A_0)\longrightarrow D^1(A\mid R',A_0)
\longrightarrow D^1(R\mid R',k)\otimes_k A_0\longrightarrow0.
\]
这意味着存在一个 $A$ 被 $A_0$ \Term{group-extension}{扩张}的 $R'$-\Term{algebra}{代数}\Term{extension}{扩张}，记作 $A'$，使得在映射
\[
D^1(A\mid R',A_0)\longrightarrow D^1(R\mid R',k)\otimes_k A_0
\]
下 $[A']\mapsto[R']\otimes1$；这就意味着\Term{deformation}{形变} $(R,\operatorname{Spec}(A))$ \Term{liftable}{可提升}为\Term{deformation}{形变} $(R',\operatorname{Spec}(A'))$。

(b) 由3.14有\Term{isomorphism}{同构}
\[
D^1(A_0\mid k,A_0)\xrightarrow\sim D^1((A_0)_x\mid k,(A_0)_x)
\xrightarrow\sim D^1(\widehat{(A_0)_x}\mid k,\widehat{(A_0)_x}),
\]
即
\[
[\mathcal M_{X_0}(k[\epsilon])]\simeq
[\mathcal M_{\operatorname{Spec}(\mathcal O_{X_0,x})}(k[\epsilon])]
\simeq[\mathcal M_{\operatorname{Spec}(\widehat{\mathcal O}_{X_0,x})}(k[\epsilon])].
\]
由2.7(a)，只需证明
\[
[\mathcal M_{X_0}]\longrightarrow[\mathcal M_{\operatorname{Spec}(\mathcal O_{X_0,x})}]
\longrightarrow[\mathcal M_{\operatorname{Spec}(\widehat{\mathcal O}_{X_0,x})}]
\]
两者都是\Term{smooth}{光滑}\Term{functor}{函子}。设 $(R,\operatorname{Spec}(A))$ 为 $X_0=\operatorname{Spec}(A_0)$ 的\Term{deformation}{形变}，$0\to k\to R'\to R\to0$ 为 $\mathcal C_\Lambda$ 中的\Term{small-extension}{小扩张}。如上，环\Term{homomorphism}{同态} $R'\to R\to A$ 给出各行\Term{exact}{正合}的\Term{commutative-diagram}{交换图}
\SourcePage{116}
\[
\resizebox{\textwidth}{!}{%
\begin{tikzcd}[ampersand replacement=\&,column sep=small,row sep=large]
D^1(A_0\mid k,A_0)\arrow[r]\arrow[d,"\sim"']\&D^1(A\mid R',A_0)\arrow[r,"p"]\arrow[d]\&D^1(R\mid R',k)\otimes_k A_0\arrow[r]\arrow[d]\&D^2(A_0\mid k,A_0)\arrow[r]\arrow[d,"\sim"]\&\cdots\\
D^1(A_0\mid k,(A_0)_x)\arrow[r]\arrow[d,equal]\&D^1(A\mid R',(A_0)_x)\arrow[r]\arrow[d]\&D^1(R\mid R',k)\otimes_k(A_0)_x\arrow[r]\arrow[d,equal]\&D^2(A_0\mid k,(A_0)_x)\arrow[r]\arrow[d,equal]\&\cdots\\
D^1((A_0)_x\mid k,(A_0)_x)\arrow[r]\&D^1(A_x\mid R',(A_0)_x)\arrow[r,"p_x"]\&D^1(R\mid R',k)\otimes_k(A_0)_x\arrow[r]\&D^2((A_0)_x\mid k,(A_0)_x)\arrow[r]\&\cdots .
\end{tikzcd}%
}
\]
其中所标的\Term{isomorphism}{同构}来自如下事实：$A_0$ 除一个\Term{closed-point}{闭点} $x$ 外在 $k$ 上\Term{smooth}{光滑}，故对所有 $i>0$ 有
\[
D^i(A_0\mid k,(A_0)_x)=D^i(A_0\mid k,A_0)_x=D^i(A_0\mid k,A_0).
\]
现设给定 $[\mathcal M_{\operatorname{Spec}((A_0)_x)}]$ 中的\Term{arrow}{箭头} $(R',\operatorname{Spec}(B))\to(R,\operatorname{Spec}(A_x))$。这就是说给定\Term{class}{类} $[B]\in D^1(A_x\mid R',(A_0)_x)$，使 $p_x([B])=[R']$。一次直接的图追踪表明，存在 $[A']\in D^1(A\mid R',A_0)$，满足 $p([A'])=[R']$ 且 $A'_x=B$；即 $(R',\operatorname{Spec}(A'))$ 是 $\operatorname{Spec}(A_0)$ 到 $R'$ 的\Term{deformation}{形变}，使 $\operatorname{Spec}(A')_x=\operatorname{Spec}(B)$。这就证明 $[\mathcal M_{X_0}]\to[\mathcal M_{\operatorname{Spec}(\mathcal O_{X_0,x})}]$ \Term{smooth}{光滑}。\Term{functor}{函子} $\mathcal M_{\operatorname{Spec}((A_0)_x)}\to\mathcal M_{\operatorname{Spec}(\widehat{(A_0)_x})}$ 的\Term{smooth}{光滑}性由类似的图及下列事实得到：对所有 $i>0$，
\[
D^i((A_0)_x\mid k,\widehat{(A_0)_x})
=\widehat{D^i((A_0)_x\mid k,(A_0)_x)}
=D^i(\widehat{(A_0)_x}\mid k,\widehat{(A_0)_x}),
\]
因为 $x$ 是唯一可能的非\Term{smooth}{光滑}点。
\end{proof}

\begin{SGAresult}{\Term{theorem}{定理}}{4.11}
设 $X_0$ 为 $k$ 上的（\Term{separated}{分离}）\Term{finite-type}{有限型}\Term{scheme}{概形}，在 $X_0$ 的\Term{finite}{有限}\Term{closed-subset}{闭集} $X_0^{\mathrm{sing}}=\{x_1,\ldots,x_r\}$ 以外\Term{smooth}{光滑}。设 $U_0^1,\ldots,U_0^r$ 分别为 $x_1,\ldots,x_r$ 的\Term{affine}{仿射}\Term{open-neighborhood}{开邻域}，且 $U_0^i\cap X_0^{\mathrm{sing}}=\{x_i\}$。若 $H^2(X_0,\mathcal D^0_{X_0})=0$，则自然\Term{functor}{函子}
\[
\mathcal M_{X_0}\longrightarrow\mathcal M_{U_0^1}\times\cdots\times\mathcal M_{U_0^r}
\]
\Term{smooth}{光滑}。
\end{SGAresult}
\begin{proof}
选择 $U_0^{r+1},\ldots,U_0^N$，使对所有 $j>r$ 有 $U_0^j\cap X_0^{\mathrm{sing}}=\varnothing$，即可将 $U_0^1,\ldots,U_0^r$ 补成 $X_0$ 的一个\Term{affine-open-covering}{仿射开覆盖}。设 $R'\to R$ 是 $\mathcal C_\Lambda$ 中的\Term{small-extension}{小扩张}，$(R',U'^1),\ldots,(R',U'^r)$ 是 $U_0^1,\ldots,U_0^r$ 到 $R'$ 的\Term{deformation}{形变}，满足对 $1\leq i\leq r$ 有 $U'^i\otimes_{R'}R=X|_{U_0^i}$，其中 $(R,X)$ 为 $X_0$ 到 $R$ 的某个\Term{deformation}{形变}。
\SourcePage{117}
须证明存在 $X_0$ 到 $R'$ 的\Term{deformation}{形变} $(R',X')$，使 $X'\otimes_{R'}R=X$，且对 $1\leq i\leq r$ 有 $X'|_{U_0^i}=U'^i$。因 $j>r$ 时 $U_0^j$ 在 $k$ 上\Term{smooth}{光滑}，由4.9(a)，$X|_{U_0^j}$ \Term{liftable}{可提升}为\Term{deformation}{形变} $(R',U'^j)$，即存在 $U_0^j$ 到 $R'$ 的\Term{deformation}{形变} $(R',U'^j)$，对所有 $j>r$ 满足 $U'^j\otimes_{R'}R=X|_{U_0^j}$。对于任意一对 $i\ne j$，$U_0^i\cap U_0^j$ 为\Term{affine}{仿射}且在 $k$ 上\Term{smooth}{光滑}，故得到\Term{isomorphism}{同构}
\[
\tau_{ij}:U'^i|_{U_0^i\cap U_0^j}\xrightarrow\sim U'^j|_{U_0^i\cap U_0^j},
\]
使 $\tau_{ij}\otimes_{R'}R=\operatorname{id}_{X|_{U_0^i\cap U_0^j}}$，进而由
\[
U_0^i\cap U_0^j\cap U_0^k\longmapsto\tau_{ij}\tau_{jk}\tau_{ik}^{-1}
\]
定义 $H^2(X_0,\mathcal D^0_{X_0})$ 中的一个\Term{cohomology-class}{上同调类}。由假设 $H^2(X_0,\mathcal D^0_{X_0})=0$，将这些\Term{object}{对象}\Term{gluing}{粘合}成\Term{global-deformation}{整体形变}的\Term{obstruction}{障碍}消失，故得到 $X_0$ 到 $R'$ 的\Term{deformation}{形变} $(R',X')$，满足 $X'\otimes_{R'}R=X$ 及 $X'|_{U_0^j}\simeq U'^j$。证明完成。
\end{proof}

\textbf{注记 4.12。} 上一证明假定 $X_0$ 为\Term{separated}{分离}，这个假设其实不必要。事实上，固定 $X_0$ 的\Term{deformation}{形变} $(R,X)$ 及 $U_0^j$ 的\Term{deformation}{形变} $(R',U'^j)$，则 $X$ 的\Term{open-subset}{开集} $V$ 到 $R'$ 的、与已给\Term{deformation}{形变} $(R',U'^j)$ 相容的\Term{deformation}{形变} $(R',V')$ 形成一个“\Term{gerbe}{凝集}”（gerbe，按 Giraud 的术语），其结构\Term{sheaf}{层}为 $\mathcal D^0_{X_0}$。所以这个“\Term{gerbe}{凝集}”存在\Term{section}{截面}（即存在与给定数据 $(R',U'^j)$ 相容的 $X$ 的\Term{global-deformation}{整体形变} $(R',X')$）的\Term{obstruction}{障碍}位于 $H^2(X_0,\mathcal D^0_{X_0})$ 中，而后者由假设为零。

\begin{SGAresult}{推论}{4.13}
设 $X_0$ 为 $k$ 上的\Term{finite-type}{有限型}\Term{scheme}{概形}，在 $X_0$ 的\Term{finite}{有限}\Term{subset}{子集} $X_0^{\mathrm{sing}}=\{x_1,\ldots,x_r\}$ 以外\Term{smooth}{光滑}，且满足 $H^2(X_0,\mathcal D^0_{X_0})=0$。则：

(a) \Term{functor}{函子}
\[
\mathcal M_{X_0}\longrightarrow
\mathcal M_{\operatorname{Spec}(\widehat{\mathcal O}_{X_0,x_1})}\times\cdots\times
\mathcal M_{\operatorname{Spec}(\widehat{\mathcal O}_{X_0,x_r})}
\]
\Term{smooth}{光滑}。特别地，
\[
\dim\mathcal M_{X_0}=\dim H^1(X_0,\mathcal D^0_{X_0})
+\sum_{c=1}^r\dim\mathcal M_{\operatorname{Spec}(\widehat{\mathcal O}_{X_0,x_c})}.
\]

(b) 若 $X_0$ 还\Term{local-scheme}{局部}是\Term{complete-intersection}{完全交}，则 $\mathcal M_{X_0}$ \Term{smooth}{光滑}。
\end{SGAresult}
\SourcePage{118}
\begin{proof}
(a) 第一个断言由4.10及4.9(b)得到。现设 $(R,\mathfrak X)$ 及 $(R_1,\mathfrak X_1),\ldots,(R_r,\mathfrak X_r)$ 分别为 $X_0$ 及 $\operatorname{Spec}(\mathcal O_{X_0,x_1}),\ldots,\operatorname{Spec}(\mathcal O_{X_0,x_r})$ 的\Term{formal-versal-deformation}{形式万有形变}。于是得到\Term{commutative-diagram}{交换图}
\[
\begin{tikzcd}[column sep=large]
h_R\arrow[r]\arrow[d]&h_{R_1\widehat\otimes_\Lambda\cdots\widehat\otimes_\Lambda R_r}\arrow[d]\\
\mathcal M_{X_0}\arrow[r]&
\mathcal M_{\operatorname{Spec}(\widehat{\mathcal O}_{X_0,x_1})}\times\cdots\times
\mathcal M_{\operatorname{Spec}(\widehat{\mathcal O}_{X_0,x_r})}.
\end{tikzcd}
\]
复合\Term{functor}{函子}
\[
h_R\longrightarrow\mathcal M_{X_0}\longrightarrow
\mathcal M_{\operatorname{Spec}(\widehat{\mathcal O}_{X_0,x_1})}\times\cdots\times
\mathcal M_{\operatorname{Spec}(\widehat{\mathcal O}_{X_0,x_r})}
\]
为\Term{smooth}{光滑}，而
\[
h_{R_1\widehat\otimes_\Lambda\cdots\widehat\otimes_\Lambda R_r}\longrightarrow
\mathcal M_{\operatorname{Spec}(\widehat{\mathcal O}_{X_0,x_1})}\times\cdots\times
\mathcal M_{\operatorname{Spec}(\widehat{\mathcal O}_{X_0,x_r})}
\]
为\Term{minimally-smooth}{极小光滑}，故由2.7(a)，
\[
h_R\longrightarrow h_{R_1\widehat\otimes_\Lambda\cdots\widehat\otimes_\Lambda R_r}
\]
\Term{smooth}{光滑}，即 $R$ 是 $R_1\widehat\otimes_\Lambda\cdots\widehat\otimes_\Lambda R_r$ 上的\Term{formal-power-series-ring}{形式幂级数环}。由4.4，下列序列\Term{exact}{正合}：
\[
\begin{aligned}
0\longrightarrow H^1(X_0,\mathcal D^0_{X_0})\longrightarrow
[\mathcal M_{X_0}(k[\epsilon])]\longrightarrow{}&
[\mathcal M_{\operatorname{Spec}(\widehat{\mathcal O}_{X_0,x_1})}(k[\epsilon])]\times\cdots\\
&\cdots\times[\mathcal M_{\operatorname{Spec}(\widehat{\mathcal O}_{X_0,x_r})}(k[\epsilon])],
\end{aligned}
\]
即
\[
\begin{aligned}
0\longrightarrow H^1(X_0,\mathcal D^0_{X_0})\longrightarrow{}
&\operatorname{Hom}_{\Lambda\text{-alg}}(R,k[\epsilon])\\
\longrightarrow{}&\operatorname{Hom}_{\Lambda\text{-alg}}(R_1\widehat\otimes_\Lambda\cdots\widehat\otimes_\Lambda R_r,k[\epsilon])
\end{aligned}
\]
\Term{exact}{正合}。因此 $R$ 是 $R_1\widehat\otimes_\Lambda\cdots\widehat\otimes_\Lambda R_r$ 上的 $d$ 个变量的\Term{formal-power-series-ring}{形式幂级数环}，其中 $d=\dim_k H^1(X_0,\mathcal D^0_{X_0})$。所以
\[
\begin{aligned}
\dim\mathcal M_{X_0}
&=\dim(k\widehat\otimes_\Lambda R)
=d+\dim\bigl(k\widehat\otimes_\Lambda(R_1\widehat\otimes_\Lambda\cdots\widehat\otimes_\Lambda R_r)\bigr)\\
&=d+\dim\bigl((k\widehat\otimes_\Lambda R_1)\widehat\otimes_k\cdots\widehat\otimes_k(k\widehat\otimes_\Lambda R_r)\bigr)\\
&=d+\sum_{i=1}^r\dim(k\widehat\otimes_\Lambda R_i)
=d+\sum_{i=1}^r\dim\mathcal M_{\operatorname{Spec}(\widehat{\mathcal O}_{X_0,x_i})}.
\end{aligned}
\]
\SourcePage{119}
(b) 由4.10及4.9(a)得到。
\end{proof}

\textbf{4.14。} 本节最后在一种简单情形中显式构造\Term{formal-versal-deformation}{形式万有形变}。事实上，若 $X_0$ 为\Term{affine}{仿射}且为\Term{complete-intersection}{完全交}，则显式写出它非常容易。设
\[
A_0=k[x_1,\ldots,x_n]/(f_1,\ldots,f_r)\qquad S/J.
\]
\SGAsourcegap{原书此式在商式 $k[x_1,\ldots,x_n]/(f_1,\ldots,f_r)$ 与 $S/J$ 之间未印连接\Term{signature}{符号}。}
其中 $S=k[x_1,\ldots,x_n]$，$J=(f_1,\ldots,f_r)$，$f_1,\ldots,f_r$ 是一个 $S$-\Term{regular-sequence}{正则序列}。于是得到\Term{exact-sequence}{正合列}
\[
\cdots\longrightarrow\operatorname{Der}_k(S,A_0)\longrightarrow
\operatorname{Hom}_{A_0}(J/J^2,A_0)\longrightarrow D^1(A_0\mid k,A_0)\longrightarrow0.
\]
注意，由于 $f_1,\ldots,f_r$ 为 $S$-\Term{regular-sequence}{正则序列}，$J/J^2$ 是由 $f_1\pmod{J^2},\ldots,f_r\pmod{J^2}$ 生成的秩 $r$ 的\Term{free-module}{自由 $A_0$-模}。现假设 $D^1(A_0\mid k,A_0)$ 是 $k$ 上的\Term{finite-dimensional}{有限维}\Term{vector-space}{向量空间}，例如 $A$ 只有孤立的非\Term{smooth}{光滑}点。令 $d=\dim_kD^1(A_0\mid k,A_0)$，选择 $p_j:J/J^2\to A_0$（$j=1,2,\ldots,d$），使其形成 $D^1(A_0\mid k,A_0)$ 的一组 $k$-基。再选择 $P_{jk}\in\Lambda[x_1,\ldots,x_n]$，使 $p_j(f_i)=P_{ji}\pmod{\mathfrak m_\Lambda+J}$，其中 $\mathfrak m_\Lambda$ 为 $\Lambda$ 的\Term{maximal-ideal}{极大理想}。考虑
\[
\begin{gathered}
R=\Lambda[[t_1,\ldots,t_d]],\\
B=R[x_1,\ldots,x_n]/(G_1,\ldots,G_r),
\end{gathered}
\]
其中
\[
G_j=F_j-\sum_{i=1}^d t_iP_{ji},
\]
而 $F_j$ 是将 $f_j$ 的系数从 $k$ \Term{lifting}{提升}到 $\Lambda$ 所得的系数属于 $\Lambda$ 的\Term{polynomial}{多项式}。因为 $R[x_1,\ldots,x_n]$ 为 $R$-\Term{flat}{平坦}，$G_1,\ldots,G_r$ 模 $\mathfrak m_R$ 后（成为 $f_1,\ldots,f_r$）是 $k\otimes_R R[x_1,\ldots,x_n]$ 中的\Term{regular-sequence}{正则序列}，所以 $B$ 为 $R$-\Term{flat}{平坦}，且 $k\otimes_R B=A_0$。此外，由我们对 $B$ 的构造，
\[
\operatorname{Hom}_{\Lambda\text{-alg}}(R,k[\epsilon])\longrightarrow
D^1(A_0\mid k,A_0)=[\mathcal M_{\operatorname{Spec}(A_0)}(k[\epsilon])]
\]
为\Term{bijection}{双射}。因此，若令
\SourcePage{120}
\[
\operatorname{Spf}(B)=\varinjlim_\nu\operatorname{Spec}(R/\mathfrak m^{\nu+1}\otimes_R B),
\]
则 $(R,\operatorname{Spf}(B))$ 为 $\widehat{\mathcal M}_{\operatorname{Spec}(A_0)}$ 中的\Term{object}{对象}；由它诱导的\Term{functor}{函子} $h_R\to\mathcal M_{\operatorname{Spec}(A_0)}$ 由1.6为\Term{smooth}{光滑}，因为 $R$ 是 $\Lambda$ 上的\Term{formal-power-series-ring}{形式幂级数环}。换言之，$(R,\operatorname{Spf}(B))$ 是 $\mathcal M_{\operatorname{Spec}(A_0)}$ 的\Term{formal-versal-object}{形式万有对象}。

由4.10(b)，若令
\[
B=R[[x_1,\ldots,x_n]]/(G_1,\ldots,G_r),
\]
则 $(R,\operatorname{Spf}(B))$ 是 $\operatorname{Spec}(A_0)$ 的\Term{formal-versal-deformation}{形式万有形变}，其中
\[
A_0=k[[x_1,\ldots,x_n]]/(f_1,\ldots,f_r).
\]

\textbf{注记 4.15。} 设 $X_0$ 是仅有孤立非\Term{smooth}{光滑}点的 $k$ 上的\Term{finite-type}{有限型}\Term{scheme}{概形}，$(R,\mathfrak X)$ 是 $X_0$ 的\Term{formal-versal-deformation}{形式万有形变}。上述构造表明，若 $X_0$ 为\Term{affine}{仿射}且\Term{local-scheme}{局部}为\Term{complete-intersection}{完全交}，则 $\mathfrak X$ 是某个\Term{finite-type}{有限型} $R$-\Term{scheme}{概形}的\Term{formalization}{形式化}。这一事实容易推广到 $X_0$ 不一定\Term{affine}{仿射}、但 $H^2(X_0,\mathcal D^0_{X_0})=0$ 的情形。

\textbf{注记 4.16。} \Term{relative-cotangent-complex}{相对余切复形}的整体理论（此处该理论归结为\Term{length}{长度}1的\Term{truncation}{截断}\Term{complex}{复形}，但这对\Term{deformation-theory}{形变理论}的许多用途已足够）使得某些结果，例如4.10，可以推广并简化。具体说，设 $S$ 为\Term{scheme}{概形}，$S_0\subset S'\subset S$ 是 $S$ 的两个\Term{subscheme}{子概形}，由满足 $JI=0$ 的\Term{ideal}{理想} $J\supset I$ 定义；$X'$ 是 $S'$ 上的\Term{flat}{平坦}\Term{scheme}{概形}，$L_\bullet=L_\bullet^{X_0/S_0}$ 是 $X_0$ 在 $S_0$ 上的\Term{relative-cotangent-complex}{相对余切复形}，为\Term{derived-category}{导出范畴} $D(X_0,\mathcal O_X)$ 的一个\Term{object}{对象}。考察所有延拓 $X'/S'$ 的 $S$ 上的\Term{flat}{平坦}\Term{scheme}{概形} $X$ 所成的\Term{category}{范畴} $\mathcal F$，即 $X$ 带有一个 $S'$-\Term{isomorphism}{同构} $X\times_S S'\simeq X'$。则：

(a) 对 $\mathcal F$ 的一个\Term{object}{对象} $X'$，其\Term{automorphism}{自同构}群\Term{canonical}{典范}\Term{isomorphism}{同构}于 $\operatorname{Ext}^0(X_0;L_\bullet,I_{X_0})$（其中 $I_{X_0}$ 是将 $I$ 看作 $S_0$ 上的\Term{module}{模}后，在 $X_0$ 上的\Term{inverse-image}{逆像}）。

\SourcePage{121}
(b) $\mathcal F$ 的\Term{object}{对象}的\Term{isomorphism-class}{同构类}之集或为空，或为群 $\operatorname{Ext}^1(X_0;L_\bullet,I_{X_0})$ 下的\Term{torsor}{挠子}（torsor，即\Term{principal-homogeneous-space}{主齐性空间}）。

(c) 可定义一个\Term{canonical}{典范}元素 $c\in\operatorname{Ext}^2(X_0;L_\bullet,I_{X_0})$，其消失为 $\mathcal F$ \Term{nonempty}{非空}的充要条件。

为研究 $\operatorname{Ext}^i(X_0;L_\bullet,I_{X_0})$，常常方便于写成
\[
\operatorname{Ext}^i(X_0;L_\bullet,I_{X_0})
=H^i\bigl(X_0,R\mathcal Hom(L_\bullet,I_{X_0})\bigr),
\]
并注意\Term{complex}{复形} $R\mathcal Hom(L_\bullet,I_{X_0})$ 的\Term{cohomology-sheaf}{上同调层}是第3节引入的\Term{sheaf}{层} $\mathcal D^i_{X_0}(I_{X_0})$。当 $i=0,1,2$ 时，这些\Term{sheaf}{层}分别可视为 $\mathcal F$ 的任一\Term{object}{对象} $X$ 的\Term{automorphism}{自同构}\Term{sheaf}{层}（意指在 $X'$ 上诱导恒等的\Term{automorphism}{自同构}）、\Term{flat}{平坦}\Term{deformation}{形变}问题的\Term{local-scheme}{局部}解的\Term{isomorphism-class}{同构类}的不确定性\Term{sheaf}{层}、以及\Term{local-scheme}{局部}解存在的\Term{obstruction}{障碍}\Term{sheaf}{层}。另一方面，标准\Term{spectral-sequence}{谱序列}给出
\[
\operatorname{Ext}^*(X_0;L_\bullet,I_{X_0})\Longleftarrow
E_2^{p,q}=H^p\bigl(X_0,\mathcal D^q_{X_0}(I_{X_0})\bigr),
\]
从而得到熟知的\Term{isomorphism}{同构}
\[
\begin{aligned}
\operatorname{Aut}(X/X')&\simeq H^0(X_0,\mathcal D^0_{X_0}(I_{X_0}))\\
&\simeq H^0\bigl(X_0,\mathcal Hom(\Omega^1_{X_0/S_0},I_{X_0})\bigr)
=\operatorname{Hom}(\Omega^1_{X_0/S_0},I_{X_0}),
\end{aligned}
\]
以及低次\Term{exact-sequence}{正合列}
\[
0\longrightarrow H^1(X_0,\mathcal D^0_{X_0})\longrightarrow\operatorname{Ext}^1
\longrightarrow H^0(X_0,\mathcal D^1_{X_0})\longrightarrow H^2(X_0,\mathcal D^0_{X_0})
\longrightarrow\operatorname{Ext}^2.
\]
特别地，这表明：若 $H^2(X_0,\mathcal D^0_{X_0})=0$，且存在 $X'/S'$ 的\Term{local-scheme}{局部}\Term{flat}{平坦}\Term{deformation}{形变} $X$，则存在与给定\Term{local-scheme}{局部}\Term{deformation}{形变}相对应的\Term{flat}{平坦}\Term{global-deformation}{整体形变}。更准确且更\Term{generic}{一般}地说，给定\Term{local-scheme}{局部}\Term{deformation}{形变}（即“\Term{local-scheme}{局部}\Term{flat}{平坦}\Term{deformation}{形变}的\Term{isomorphism-class}{同构类}\Term{sheaf}{层}”的一个\Term{section}{截面}），4.10的论证表明，寻找与之相容的\Term{flat}{平坦}\Term{global-deformation}{整体形变}的\Term{obstruction}{障碍}属于 $H^2(X_0,\mathcal D^0_{X_0})$；故若 $H^2(X_0,\mathcal D^0_{X_0})=0$，这样的\Term{flat}{平坦}\Term{global-deformation}{整体形变}存在：这本质上就是4.10。

\SourcePage{122}
另一方面，\Term{spectral-sequence}{谱序列}表明，只要满足
\[
H^0(X_0,\mathcal D^2_{X_0})=0,\qquad
H^1(X_0,\mathcal D^1_{X_0})=0,\qquad
H^2(X_0,\mathcal D^0_{X_0})=0,\tag{*}
\]
则 $\operatorname{Ext}^2$ 为零。若 $X_0$ 在 $S_0$ 上为\Term{relative-complete-intersection}{相对完全交}，第一个关系总成立（这时对 $X_0$ 上的任意\Term{module}{模} $M$，有 $\mathcal D^2_{X_0}=0$）；若 $\mathcal D^1_{X_0}$ 的\Term{support}{支集}为\Term{discrete}{离散}的，则第二个关系成立，例如 $X_0$ 除去 $X_0$ 的某个\Term{discrete}{离散}\Term{subset}{子集}后在 $S_0$ 上\Term{smooth}{光滑}时便是如此。上述三个关系 (*) 均成立时，由上面的(c)，$\mathcal F$ \Term{nonempty}{非空}；这推广了4.12(b)。

从更几何的角度，无须使用\Term{relative-cotangent-complex}{相对余切复形}的整体理论（只需将\Term{functor}{函子} $\mathcal D^i_{X_0}$ 整体化），便可对将 $X'$ \Term{flat}{平坦}地\Term{deformation}{形变}为 $X$ 的三个相继\Term{obstruction}{障碍}给出简单的几何解释；它们分别位于 (*) 左侧的三个项中。具体地，由\Term{local-scheme}{局部}理论，第一个\Term{obstruction}{障碍} $c_1\in H^0(X_0,\mathcal D^2_{X_0})$ 可定义为寻找\Term{local-scheme}{局部}解的\Term{obstruction}{障碍}。此\Term{obstruction}{障碍}消失时，由\Term{local-scheme}{局部}理论，问题的\Term{local-scheme}{局部}解的\Term{isomorphism-class}{同构类}\Term{sheaf}{层} $\mathcal F$ \Term{canonical}{典范}地成为 $\mathcal D^1_{X_0}$ 下的\Term{torsor}{挠子}；位于 $H^1(X_0,\mathcal D^1_{X_0})$ 中的第二个\Term{obstruction}{障碍} $c_2$，就是寻找这个\Term{torsor}{挠子}的一个\Term{section}{截面}的\Term{obstruction}{障碍}，即寻找“\Term{local-scheme}{局部}解的\Term{isomorphism-class}{同构类}的一个相容系统”的\Term{obstruction}{障碍}。若第二个\Term{obstruction}{障碍}也消失，则任取 $\mathcal F$ 的一个\Term{section}{截面}，如上所述，得到一个 $H^2(X_0,\mathcal D^0_{X_0})$ 中的\Term{obstruction}{障碍} $c_3$，它是寻找与给定的\Term{local-scheme}{局部}解（的\Term{isomorphism-class}{同构类}）系统相对应的整体解的\Term{obstruction}{障碍}。利用 $\mathcal F$ 的\Term{section}{截面}确定到加上 $\mathcal D^1_{X_0}$ 的一个\Term{section}{截面}这一事实，得到映射
\[
d:H^0(X_0,\mathcal D^1_{X_0})\longrightarrow H^2(X_0,\mathcal D^0_{X_0}),
\]
且\Term{deformation}{形变}问题存在解，当且仅当 $c_1=0$、$c_2=0$，且 $c_3$ 作为 $\operatorname{Coker}d$ 中的元素为零（每个 $c_i$ 都在此前的各个\Term{obstruction}{障碍}为零时才有定义）。当然，群 $H^0(X_0,\mathcal D^2_{X_0})$、$H^1(X_0,\mathcal D^1_{X_0})$ 以及
\SourcePage{123}
$H^2(X_0,\mathcal D^0_{X_0})/\operatorname{Im}d$ 正是上述\Term{spectral-sequence}{谱序列}的起始项 $E_2^{0,2}$、$E_2^{1,1}$ 和 $E_3^{2,0}$，而 $d$ 就是\Term{differential}{微分}（相差一个\Term{signature}{符号}？）
\[
d_2:E_2^{0,1}\longrightarrow E_2^{2,0}.
\]
